\bookchapter{บทที่ 5: ทำให้โทรศัพท์เข้าใจคุณมากขึ้น — การตั้งค่าพื้นฐาน}

\chapterintro{ใช้เวลาตั้งค่าแค่ไม่กี่นาที แต่ประสบการณ์การใช้งานดีขึ้นมาก การตั้งค่าเหล่านี้เหมือน “จูนเสียง” ให้โทรศัพท์เข้ากับนิสัยการใช้งานของคุณ}

\booksection{หนึ่ง เสียงและการสั่น}

ตั้งเสียงเรียกเข้าและเสียงแจ้งเตือนที่คุณชอบ ใช้โหมดเงียบ (ประชุม) กับโหมดห้ามรบกวน (นอน) ให้เป็นเวลา โลกจะเงียบสงบทันที

\begin{itemize}
  \item \textbf{เสียงเรียกเข้า}: เลือกอันที่ได้ยินชัดและไม่น่าอาย
  \item \textbf{โหมดสั่น}: ตอนเงียบพึ่งการสั่นเตือน
  \item \textbf{โหมดห้ามรบกวน}: ตั้งข้อยกเว้นได้ เช่นอนุญาตสายจากคนในสมุดโทรศัพท์
  \item \textbf{แยกระดับเสียง}: เสียงสื่อ เสียงเรียกเข้า เสียงปลุก ปรับแยกกันได้
\end{itemize}

\booksection{สอง หน้าจอและความสว่าง}

\begin{itemize}
  \item \textbf{ความสว่างอัตโนมัติ}: ให้โทรศัพท์ปรับตามแสงรอบตัวเอง สบายใจ
  \item \textbf{โหมดปกป้องตา/โหมดกลางคืน}: กรองแสงสีน้ำเงิน ดูมือถือตอนกลางคืนสบายตาขึ้น
  \item \textbf{โหมดมืด}: ประหยัดไฟ (เห็นผลชัดกับจอ OLED) และไม่แสบตาตอนกลางคืน แอปจำนวนมากก็รองรับโหมดมืดด้วย
  \item \textbf{อัตรารีเฟรชหน้าจอ}: บางรุ่นแอนดรอยด์ปรับได้ ยิ่งสูง (เช่น 120Hz) ยิ่งปัดลื่น
  \item \textbf{ล็อกหน้าจออัตโนมัติ}: ตั้ง 15 หรือ 30 วินาที ปลอดภัยและประหยัดไฟ
  \item \textbf{ขนาดตัวอักษร}: รู้สึกตัวหนังสือเล็ก? ขยายขึ้นหน่อย
\end{itemize}

\booksection{สาม วอลเปเปอร์ ธีม และหน้าจอหลัก}

เปลี่ยนวอลเปเปอร์สวยๆ สักใบ (นิ่งหรือเคลื่อนไหวก็ได้) อารมณ์ก็เปลี่ยน โทรศัพท์แบรนด์แอนดรอยด์มักมีร้านธีม เปลี่ยนไอคอนกับสไตล์ทั้งชุดได้

\textbf{วิดเจ็ตบนหน้าจอหลัก}: วางพยากรณ์อากาศ ตารางนัดหมาย ปุ่มควบคุมเพลงไว้บนหน้าจอหลัก เข้าถึงได้ในขั้นตอนเดียว ทั้ง iOS และแอนดรอยด์รองรับ กดค้างที่ที่ว่างบนหน้าจอหลักเพื่อเพิ่ม

\booksection{สี่ การตั้งค่าเพื่อประสิทธิภาพ: ให้โทรศัพท์เป็น “สมองเสริม” ของคุณ}

ส่วนนี้คือฟังก์ชันแบบ “ตั้งแล้วไม่อยากกลับไปแบบเดิม” แต่ละอย่างช่วยเพิ่มประสิทธิภาพชีวิตประจำวันได้ชัดเจน

\booksubsection{โหมดโฟกัส / ห้ามรบกวน: ชิงความสนใจคืนจากการแจ้งเตือน}

\begin{itemize}
  \item \textbf{แนวคิดหลัก}: ไม่ใช่ "ดูมือถือน้อยลง" แต่คือ "ตอนควรโฟกัสก็โฟกัส ตอนควรติดต่อค่อยติดต่อ"
  \item \textbf{คำแนะนำการใช้งาน}:\newline : · โหมดทำงาน: อนุญาตเฉพาะผู้ติดต่อ "ติดดาว" และแอปที่จำเป็น\newline · โหมดนอน: เปิดอัตโนมัติ อนุญาตแค่นาฬิกาปลุก\newline · ออกกำลังกาย / อ่านหนังสือ / ประชุม: ปรับตามกิจกรรมปัจจุบัน
  \item \textbf{วิธีเล่นขั้นสูง}: ใช้ระบบอัตโนมัติ + ตัวกระตุ้นตามตำแหน่ง/เวลา ให้โหมดสลับเอง เช่น คำสั่งลัดของ iOS ตั้ง "พอถึงบริษัทให้เข้าโหมดทำงานอัตโนมัติ" ได้
\end{itemize}

\booksubsection{การจัดการการแจ้งเตือน: จุดแดงไม่รบกวนชีวิตอีกต่อไป}

\begin{itemize}
  \item \textbf{ปิดการแจ้งเตือนสายโฆษณา}: การแจ้งเตือน 90\% แท้จริงคุณไม่จำเป็นต้องเห็นเลย
  \item \textbf{จัดลำดับตามความสำคัญ}: ติดต่อสื่อสารกับคน (โทรศัพท์ SMS แชท) ให้ขึ้นบนสุด ส่วนข่าว ช้อปปิ้ง โซเชียล ตั้งเป็น "แจ้งเงียบ" ได้
  \item \textbf{จัดกลุ่มการแจ้งเตือน}: พับการแจ้งเตือนประเภทเดียวกันไว้ด้วยกัน กันการปลดล็อกทีเดียวโดนข้อความนับสิบท่วม
\end{itemize}

\booksubsection{ท่าทางลัด: ให้การใช้งานเร็วกว่าคนอื่นหนึ่งก้าว}

\begin{itemize}
  \item \textbf{จับภาพหน้าจอ}: ปุ่มพาวเวอร์ + ลดเสียง (iOS จอเต็ม / แอนดรอยด์ทั่วไป)
  \item \textbf{ภาพยาว / จับหน้าจอแบบเลื่อน}: จับภาพปกติเสร็จแล้วกดตัวเลือก "ภาพยาว"
  \item \textbf{สแกนคิวอาร์ปุ๊บปั๊บ}: เพิ่มไอคอน "สแกน" ไว้ในศูนย์ควบคุม/การตั้งค่าด่วน
  \item \textbf{แตะด้านหลังเครื่อง (iPhone)}: ผูก "จับภาพหน้าจอ" หรือ "ล็อกหน้าจอ" ฯลฯ ไว้ที่ [ตั้งค่า] > [การช่วยการเข้าถึง] > [สัมผัส] > [แตะด้านหลัง]
  \item \textbf{เรียกด้วยเสียง}: Siri / Xiaoai / Xiaoyi / Xiaobo ตะโกนหนึ่งคำก็ถามอากาศ ตั้งปลุกได้
\end{itemize}

\booksubsection{คีย์บอร์ด: เครื่องมือเพิ่มประสิทธิภาพที่ถูกมองข้ามที่สุด}

\begin{itemize}
  \item \textbf{เลือกคีย์บอร์ดให้ถูกตัว}:\newline : · Sogou / Baidu: คลังคำใหญ่ คำสแลงอัปเดตไว\newline · iFlytek: แม่นยำที่สุดเรื่องพิมพ์ด้วยเสียง เจ้าประโยชน์ตอนจดบันทึกในที่ประชุม\newline · Gboard (Google): เรียบง่าย ซิงก์ข้ามแพลตฟอร์มเก่ง พิมพ์ปัดได้\newline · คีย์บอร์ดในตัวของ iOS: ดูแลความเป็นส่วนตัวดีสุด และเชื่อมกับระบบลึก
  \item \textbf{สวิตช์ที่ควรเปิด}: แก้พิมพ์ผิดอัตโนมัติ คำทายต่อ ซิงก์คลังคำคลาวด์ ปัดขึ้นพิมพ์ตัวเลขสัญลักษณ์
  \item \textbf{เตือนเรื่องความเป็นส่วนตัว}: ใช้ฟังก์ชัน "ทายทั้งประโยค" หรือ "แนะนำเฉพาะบุคคล" อย่างระมัดระวัง อาจอัปโหลดสิ่งที่คุณพิมพ์
\end{itemize}

\booksubsection{หน้าจอหลักกับมัลติทาสก์: ให้เวิร์กโฟลว์ลื่นไหล}

\begin{itemize}
  \item \textbf{แบ่งหน้าจอแอป / ภาพทับ (PiP)}: ดูวิดีโอไปพลางตอบแชท เล่นเว็บไปพลางจดโน้ต
  \item \textbf{หน้าต่างเล็ก / หน้าต่างลอย}: WeChat / Douyin รองรับ "โหมดหน้าต่างเล็ก" ตอบเสร็จก็หุบ
  \item \textbf{ปักหมุดแอปที่ใช้บ่อย}: แถบ Dock (ด้านล่างจอ) วางแค่ 4 ตัวที่ใช้บ่อยสุด ลดเวลาค้นหา
\end{itemize}

\begin{tcbinfo}{เช็กลิสต์ "วันแรกของการใช้งาน"}
ได้เครื่องใหม่หรืออัปเดตเวอร์ชันใหญ่แล้ว แนะนำใช้เวลา 10 นาทีไล่ตั้งค่าพวกนี้รอบหนึ่ง:\newline (1) รหัสล็อกหน้าจอ + ไบโอเมตริกซ์ (2) จัดกลุ่มการแจ้งเตือน + ปิดพุชโฆษณา (3) ขนาดตัวอักษร / อัตรารีเฟรช / โหมดมืด\newline (4) คีย์บอร์ด / เบราว์เซอร์ / อีเมลเริ่มต้น (5) ผู้ติดต่อฉุกเฉิน / บัตรข้อมูลการแพทย์ (6) การตั้งค่าสำรองข้อมูล

\end{tcbinfo}

\booksection{ห้า การช่วยการเข้าถึง: พลังของ “เทคโนโลยีเพื่อมนุษย์”}

อย่ากลัวชื่อ “การช่วยการเข้าถึง” ฟังก์ชันพวกนี้ใช้ได้กับทุกคน แรกเริ่มออกแบบเพื่อผู้พิการ แต่คนทั่วไปก็ใช้ประโยชน์ได้มาก — มือบาดเจ็บชั่วคราวแล้วต้องอุ้มลูก หูฟังพังบนรถไฟใต้ดินเลยไม่ได้ยิน ใส่แว่นสายตายาวแล้วอ่านตัวเล็กไม่ออก… ล้วนช่วยได้ทั้งนั้น ท้ายที่สุดการออกแบบเพื่อการเข้าถึงสะท้อนแนวคิด\textbf{เทคโนโลยีเพื่อมนุษย์}: ให้ทุกคนใช้ชีวิตดิจิทัลได้อย่างเท่าเทียมกัน

\booksubsection{“ดวงตา” ของผู้ใช้ทาสายตา — คู่มือฉบับเต็มของโปรแกรมอ่านหน้าจอ}

โปรแกรมอ่านหน้าจอคือเครื่องมือหลักของผู้เสียการมองเห็น มันจะอ่าน\textbf{ทุกเนื้อหาที่มองเห็นบนหน้าจอออกเสียง}ให้ฟัง ผู้ใช้ควบคุมโทรศัพท์ด้วยท่าทางและเสียง ต่อไปแนะนำวิธีใช้ VoiceOver ของ iOS กับ TalkBack ของแอนดรอยด์

\begin{tcbinfo}{iOS VoiceOver (โปรแกรมอ่านหน้าจอของแอปเปิล)}
\textbf{เส้นทางเปิด}: [ตั้งค่า] > [การช่วยการเข้าถึง] > [VoiceOver] เปิดสวิตช์ หรือกดปุ่มด้านข้าง (หรือปุ่มโฮม) สามครั้งเพื่อเปิดปิดด่วน

\textbf{ท่าทางพื้นฐาน}:

\begin{enumerate}
  \item : ปัดนิ้วเดียว: เลื่อนดูเนื้อหาบนจอทีละองค์ประกอบ
  \item ปัดสองนิ้ว: เลื่อนทั้งหน้าขึ้นลง/ซ้ายขวา
  \item แตะสองครั้งด้วยนิ้วเดียว: เปิดใช้องค์ประกอบที่โฟกัสอยู่ (เท่ากับการ "แตะ")
  \item ปัดสามนิ้ว: สลับหน้า
  \item แตะสามครั้งด้วยนิ้วเดียว: ซูมหรือเลือก (ต่างจากการแตะสามครั้งทั่วไป)
\end{enumerate}

\textbf{การตั้งค่าเฉพาะบุคคล}: ที่ [VoiceOver] > [เสียงพูด] ปรับได้ว่า\textbf{ความเร็วในการอ่าน}(ค่าเริ่มต้นราว 50\% ปรับ 70-80\% ฟังลื่นกว่า) และ\textbf{ระดับเสียง}, เปลี่ยนชุดเสียงได้ iOS มี “เสียงขั้นสูง” และ “ตัวควบคุมเสียง” ในตัว อันแรกช่วยลดเสียงหายใจและการหยุดชะงัก ฟังธรรมชาติกว่า

\end{tcbinfo}

\begin{tcbinfo}{แอนดรอยด์ TalkBack}
\textbf{เส้นทางเปิด}: [ตั้งค่า] > [การช่วยการเข้าถึง] > [TalkBack] เปิดสวิตช์ หรือกดปุ่มเพิ่ม-ลดเสียงค้าง 3 วินาทีเพื่อเปิดด่วน (บางแบรนด์รองรับ)

\textbf{ท่าทางพื้นฐาน}:

\begin{enumerate}
  \item : ปัดนิ้วเดียวซ้ายขวา: สลับองค์ประกอบที่โฟกัส
  \item แตะสองครั้งนิ้วเดียว: เปิดใช้องค์ประกอบ
  \item ปัดสองนิ้วขึ้นลง/ซ้ายขวา: เลื่อนหน้าจอ
  \item แตะครั้งเดียวด้วยสามนิ้ว: เข้าโหมดอ่าน (อ่านต่อเนื่อง)
  \item แตะครั้งเดียวด้วยสี่นิ้ว: หยุด/อ่านต่อ
\end{enumerate}

\textbf{การตั้งค่าเฉพาะบุคคล}: ที่ [TalkBack] > [การตั้งค่าการแปลงข้อความเป็นคำพูด] ปรับได้ว่า\textbf{ความเร็วเสียงพูด},\textbf{ระดับเสียง}, และเปลี่ยน\textbf{เอนจินเสียงพูด}(แนะนำติดตั้ง Google Text-to-Speech เสียงจะธรรมชาติกว่า)

\end{tcbinfo}

\booksubsection{ทำตัวหนังสือเล็กให้ใหญ่ ภาพเบลอให้ชัด — การปรับด้านการมองเห็น}

\booksubsection{ท่าทางซูม}

\begin{itemize}
  \item \textbf{iOS}: [ตั้งค่า] > [การช่วยการเข้าถึง] > [ซูม] เปิดแล้ว\textbf{แตะสามนิ้วที่หน้าจอสองครั้ง}ก็ขยายได้ ค่าเริ่มต้นซูมทั้งจอ เลือกซูมแบบหน้าต่าง (แผ่นขยายเฉพาะจุด) ก็ได้\textbf{ใช้สองนิ้วถ่าง-หุบเพื่อปรับระดับการขยาย (1.5x ถึง 15x)}แล้วแตะสองครั้งด้วยสามนิ้วเพื่อล็อกสัดส่วนซูมปัจจุบัน
  \item \textbf{แอนดรอยด์}: [ตั้งค่า] > [การช่วยการเข้าถึง] > [ขยาย] แตะสามนิ้วเพื่อขยาย แตะซ้ำเพื่อออก หรือใช้\textbf{แตะสามนิ้วสามครั้งบนหน้าจอ}เปิดการขยายชั่วคราว (สูงสุด 8 เท่า)
\end{itemize}

\booksubsection{ตัวหนังสือคอนทราสต์สูง / ตัวหนา}

\begin{itemize}
  \item \textbf{iOS}: [ตั้งค่า] > [การช่วยการเข้าถึง] > [การแสดงผลและขนาดข้อความ] เปิดสวิตช์ "ตัวหนังสือหนา" "เพิ่มคอนทราสต์" "ลดความโปร่งใส" ฯลฯ "กลับสีอัจฉริยะ" ทำให้ตัวหนังสือยังเป็นสีดำบนพื้นสว่าง กันพื้นขาวตัวขาว
  \item \textbf{แอนดรอยด์}: [ตั้งค่า] > [การแสดงผล] > [ขนาดฟอนต์] ปรับได้สี่ระดับ "เล็ก ปกติ ใหญ่ ใหญ่พิเศษ" Huawei/Xiaomi มีสวิตช์ "ตัวหนา" และ "ข้อความคอนทราสต์สูง" เพิ่มให้ด้วย
\end{itemize}

\booksubsection{ฟิลเตอร์สี (เป็นมิตรกับคนตาบอดสี)}

\begin{itemize}
  \item \textbf{iOS}: [ตั้งค่า] > [การช่วยการเข้าถึง] > [การแสดงผลและขนาดข้อความ] > [ฟิลเตอร์สี] มีให้\textbf{โหมดตาบอดสีแดง/เขียว/น้ำเงิน/เหลืองสี่แบบ}พร้อมตัวเลือก "ระดับสีเทา" และ "กลับสี" แนะนำทดสอบการมองสีก่อน (ระบบมีให้) แล้วค่อยเลือกฟิลเตอร์ที่ตรง
  \item \textbf{แอนดรอยด์}: [ตั้งค่า] > [การช่วยการเข้าถึง] > [การแก้ไขสี] มี\textbf{สามโหมดตาบอดสี}+ สวิตช์ "กลับสี"
\end{itemize}

\booksubsection{AI บรรยายภาพ: ให้ AI เป็น “ดวงตา” ของคุณ}

เมื่อโปรแกรมอ่านหน้าจอเจอรูปภาพ วิธีเดิมคือข้ามไป หรืออ่านแค่คำว่า “รูปภาพ” แต่ตอนนี้ AI สามารถ\textbf{บรรยายเนื้อหารูปภาพด้วยประโยคเดียว}ช่วยเรื่องแคปชา การอ่านรูป การอ่านเมนู ได้มากมาย

\begin{tcbinfo}{สถานะการรองรับ AI บรรยายภาพในแต่ละระบบ}
\begin{itemize}
  \item \textbf{iOS 15+}: ที่ [ตั้งค่า] > [การช่วยการเข้าถึง] > [VoiceOver] > [คำบรรยายเสียง] เปิด "การบรรยายรูปภาพ" ระบบจะเรียกใช้\textbf{โมเดล AI ในเครื่อง}สร้างคำบรรยายรูปภาพให้อัตโนมัติ (ต้อง iPhone 12 ขึ้นไป)
  \item \textbf{แอป TalkBack เวอร์ชันใหม่}: Google TalkBack รุ่นใหม่รวมฟังก์ชัน\textbf{ถามตอบเรื่องรูปภาพ}ไว้แล้ว ถามได้ว่า "ในรูปนี้มีอะไร" AI จะตอบสดๆ (ไม่ผูกติดเวอร์ชันแอนดรอยด์เคร่งครัด อัปเดตแอป TalkBack ก็ใช้ได้)
  \item \textbf{แอปภายนอก}: "Seeing AI" (iOS) "Lookout" (แอนดรอยด์) "Xiaoyi Bangbang" (Huawei) ล้วนเป็นผู้ช่วยสายตา AI ที่ดีเยี่ยม จำแนกสิ่งของ อ่านเอกสาร รู้จำใบหน้าและสีหน้าได้
\end{itemize}

\end{tcbinfo}

\begin{tcbpractice}{สถานการณ์ใช้งาน}
ลองถามรูปถ่ายถนนสักใบว่า “ในรูปมีกี่คน กำลังทำอะไร” หรือเปิดแอปสั่งอาหาร ให้ AI อ่านชื่อเมนูจากรูป —\textbf{นี่คือ “เทคโนโลยีทำให้ชีวิตเท่าเทียม” อย่างแท้จริง}.

\end{tcbpractice}

\booksubsection{“หู” ของผู้ใช้ทาสายตา — คู่มือฉบับเต็มของคำบรรยายเรียลไทม์}

คำบรรยายเรียลไทม์ (Live Captions) คือฟังก์ชัน\textbf{แปลงเสียงพูดเป็นข้อความระดับระบบ}มีประโยชน์มากกับผู้มีปัญหาการได้ยิน ทั้งในที่เสียงดัง ห้องสมุด และสถานการณ์ภาษาต่างประเทศ

\begin{tcbinfo}{การเปิดและการปรับคำบรรยายเรียลไทม์}
\textbf{iOS}: [ตั้งค่า] > [การช่วยการเข้าถึง] > [คำบรรยายเรียลไทม์] เปิดสวิตช์ ระบบจะตรวจจับเสียงทั้งหมดของอุปกรณ์ (โทร วิดีโอ พอดแคสต์ FaceTime บันทึกเสียง) แล้วโชว์คำบรรยายเหนือหน้าจออัตโนมัติ

\textbf{แอนดรอยด์ 12+}: กด\textbf{ปุ่มปรับเสียง}แตะไอคอน “คำบรรยายเรียลไทม์” ในกล่องป๊อปอัป หรือเปิดที่ [ตั้งค่า] > [การช่วยการเข้าถึง] > [คำบรรยายเรียลไทม์]

\textbf{เคล็ดลับเพิ่มความแม่นยำ}:

\begin{enumerate}
  \item : ควรใช้ใน\textbf{สภาพแวดล้อมที่เงียบ}ห่างจากแอร์และเครื่องจักรเสียงกึกก้อง
  \item ผู้พูด\textbf{ออกเสียงชัด จังหวะกำลังดี}จะได้อัตราการจำแนกสูงสุด
  \item บทสนทนาหลายคนทำให้แม่นยำลง ใช้โหมด "แสดงเฉพาะคนเดียว" ได้
  \item รองรับภาษาหลัก เช่นจีนกลางมาตรฐาน กวางตุ้ง อังกฤษ ฯลฯ ส่วน\textbf{ภาษาเล็กแม่นยำต่ำกว่า};
  \item กลางประโยคมีเพลงประกอบหรือเสียงคนพูดแทรก การจำแนกจะสะดุด
\end{enumerate}

\end{tcbinfo}

\booksubsection{สองช่องทางสำหรับการโทร: เสียงเป็นข้อความ + ข้อความเป็นเสียง}

เมื่อผู้ใช้\textbf{ทั้งไม่ได้ยินและพูดไม่ได้}(เช่นเสียงหายฉับพลัน พักฟื้นหลังผ่าตัดคอ การพูดติดอ่างรุนแรง) โทรศัพท์ยังช่วยเป็น “คนส่งสาย” ได้

\begin{tcbinfo}{การแปลบทสนทนาเรียลไทม์ (TTY / RTT)}
\textbf{iOS}: [ตั้งค่า] > [การช่วยการเข้าถึง] > [RTT/TTY] เปิดแล้วตอนโทรหน้าจอจะมีคีย์บอร์ดขึ้นมา เสียงของอีกฝ่ายจะถูกพิมพ์เป็นข้อความแบบสดๆ\textbf{ข้อความที่คุณพิมพ์ก็จะถูกอ่านออกเสียง}ให้อีกฝ่ายฟัง

\textbf{แอนดรอยด์}: [ตั้งค่า] > [การช่วยการเข้าถึง] > [การถอดเสียงเรียลไทม์] + [การแจ้งเตือนด้วยเสียง] “การถอดเสียงเรียลไทม์”แสดงคำพูดของอีกฝ่ายเป็นข้อความสดๆ “การแจ้งเตือนด้วยเสียง”เปลี่ยนสิ่งที่คุณพิมพ์เป็นเสียงเล่นให้อีกฝ่ายฟังเวลาพูดไม่สะดวก

\textbf{ความเข้ากันได้}:

\begin{itemize}
  \item \textbf{RTT (Real-Time Text)}: ทั้งสองฝ่ายต้องรองรับ ข้อความโผล่ทีละตัวระหว่างพิมพ์ ไม่มีดีเลย์ ผู้ให้บริการหลักรองรับแล้ว
  \item \textbf{TTY (เครื่องพิมพ์ดีดส่งไกล)}: ต้องมีอุปกรณ์ TTY เฉพาะ กำลังถูก RTT แทนที่ทีละนิด
  \item ฟังก์ชัน\textbf{เปลี่ยนเสียงเป็นข้อความ}ของ WeChat/QQ ก็เป็นทางเลือกหนึ่ง: กดค้างที่ข้อความเสียง "แปลงเป็นข้อความ"
\end{itemize}

\end{tcbinfo}

\booksubsection{ฟังก์ชัน “เล็กๆ แต่ใช้ได้จริง” สำหรับทุกคน}

\begin{itemize}
  \item \textbf{แตะด้านหลังเครื่อง}(iOS 14+): [ตั้งค่า] > [การช่วยการเข้าถึง] > [สัมผัส] > [แตะด้านหลัง] ผูก\textbf{แตะสองครั้ง/สามครั้ง}เพื่อทริกเกอร์การจับภาพหน้าจอ ล็อกเครื่อง เปิดศูนย์ควบคุม ฯลฯ สุดยอดเวลาสองมือถือพัสดุ
  \item \textbf{AssistiveTouch / ลูกลอย}: iOS [การช่วยการเข้าถึง] > [สัมผัส] > [AssistiveTouch] แอนดรอยด์ [การช่วยการเข้าถึง] > [การนำทางแบบลอย] ปุ่มเสมือนแทนท่าทางหรือปุ่มกายภาพ ช่วยเวลาปุ่มไม่ค่อยทำงานหรือใช้มือเดียว
  \item \textbf{Switch Control}: ควบคุมโทรศัพท์ผ่านสวิตช์ภายนอก การแตะครั้งเดียว หรือการขยับศีรษะ\textbf{ออกแบบเพื่อผู้มีความบกพร่องทางร่างกาย}iOS [การช่วยการเข้าถึง] > [Switch Control] แอนดรอยด์ [การช่วยการเข้าถึง] > [การควบคุมด้วยสวิตช์]
  \item \textbf{การจำแนกเสียง}(iOS): [การช่วยการเข้าถึง] > [การจำแนกเสียง] เปิดแล้วโทรศัพท์จะ "ฟัง" เสียงกริ่งประตู เสียงเด็กร้องไห้ สัญญาณเตือนอัคคีภัย เสียงไซเรน ฯลฯ\textbf{จำแนกได้ก็สั่น + แจ้งเตือนบนหน้าจอ}มีค่ามากสำหรับผู้มีปัญหาการได้ยินและช่วงเวลาพักผ่อน
  \item \textbf{แจ้งเตือนด้วยไฟแฟลช}: iOS [การช่วยการเข้าถึง] > [เสียงและภาพ] > [แฟลชไฟ LED เมื่อมีการแจ้งเตือน] แอนดรอยด์ [ตั้งค่า] > [การแจ้งเตือน] > [แจ้งเตือนด้วยแสงแฟลช] แฟลชจะกะพริบเมื่อมีสายเรียกเข้าหรือแจ้งเตือน\textbf{ในที่เสียงดังก็รับรู้ได้แม้ไม่มีเสียง}.
  \item \textbf{เสียงโมโน}: [ตั้งค่า] > [การช่วยการเข้าถึง] > [เสียง] > [เสียงโมโน] รวมช่องซ้ายขวาเข้าด้วยกัน ผู้ใช้ที่หูข้างหนึ่งได้ยินน้อยจะไม่เสียข้อมูลเสียงอีกฝั่ง
\end{itemize}

\begin{tcbinfo}{ปรัชญาการออกแบบเพื่อการเข้าถึงของแอปเปิล}
แอปเปิลลงทุนกับการช่วยการเข้าถึงมหาศาล อัปเดตและเพิ่มฟังก์ชันใหม่ทุกปี เว็บทางการของแอปเปิลมีหน้าสำหรับการช่วยการเข้าถึงโดยเฉพาะ แนะนำวิธีใช้แต่ละฟังก์ชันอย่างละเอียด ฟังก์ชันการช่วยการเข้าถึงของ iOS มักสมบูรณ์กว่าและลื่นกว่าแอนดรอยด์ นี่ก็เป็นเหตุผลสำคัญที่ผู้ใช้พิการจำนวนมากเลือกแอปเปิล

\end{tcbinfo}

\begin{tcbpractice}{ลองสัมผัสดูใน 30 วินาที}
หาเวลาว่างสักช่วง แล้วลองทำดังนี้:

\begin{enumerate}
  \item เปิด\textbf{VoiceOver / TalkBack}แล้วลองใช้สามนิ้วเลื่อนดู WeChatโดยไม่ใช้ตา
  \item เปิด\textbf{คำบรรยายเรียลไทม์}แล้วพูดใส่โทรศัพท์ดูอัตราการจำแนก
  \item ใช้\textbf{การแตะสามนิ้ว}ขยายข้อความสักช่วงถึง 200\% สัมผัสชีวิตประจำวันของผู้ทาสายตา
  \item ปิดฟังก์ชันทั้งหมดที่เปิดไว้ แล้วคิดดูว่า "ถ้าฉันต้องใช้ฟังก์ชันเหล่านี้ตลอดไป ฉันอยากให้มันออกแบบมาอย่างไร"
\end{enumerate}

\end{tcbpractice}
